\documentclass[11pt,a4paper]{article}

\usepackage[utf8]{inputenc}
\usepackage[T1]{fontenc}
\usepackage[american]{babel}
\usepackage{csquotes}
\usepackage[margin=2.5cm]{geometry}
\usepackage[hidelinks]{hyperref}

% APA 7th edition citation style. NOT in Tectonic's 2022 bundle —
% pulled from CTAN at build time via the `ctan_packages` attribute
% in BUILD.bazel.
\usepackage[style=apa,backend=biber]{biblatex}
\addbibresource{references.bib}

\title{Language Games and the Limits of Description}
\author{A. Reader \\ \texttt{a.reader@example.edu}}
\date{\today}

\begin{document}
\maketitle

\begin{abstract}
This short note revisits the role of language games
in late Wittgensteinian philosophy~\parencite{wittgenstein1953}
and contrasts it with the linguistic-turn reading proposed by
\textcite{rorty1979} and elaborated in subsequent secondary
literature~\parencite{kripke1982,mcdowell1994}.
\end{abstract}

\section{Introduction}

Wittgenstein's \emph{Philosophical Investigations} reframes meaning
as use rather than reference~\parencite{wittgenstein1953}. The shift
recasts large parts of the philosophy-of-language project: not as
an attempt to specify the truth-conditions of sentences, but as a
description of the practices in which sentences earn their grip.

\textcite{rorty1979} reads this as the decisive break with
representationalist epistemology, situating Wittgenstein alongside
Heidegger and Dewey as the central twentieth-century post-Cartesian.
\textcite{kripke1982} dissents, recovering a sceptical reading on
which the rule-following passages threaten the very notion of
determinate meaning. \textcite{mcdowell1994} navigates between the
two, defending a quietist response that accepts the descriptive
reorientation without conceding the sceptical conclusion.

\section{The Argument}

\dots

\printbibliography

\end{document}
